\section{Introduction}
\label{sec:intro}

When a French speaker says \emph{quatre-vingt-dix-sept}, they are not merely naming a number---they are performing implicit arithmetic: $4 \times 20 + 10 + 7 = 97$.
This vigesimal (base-20) system, used for numbers 70--99 in standard French, stands in sharp contrast to the regular decimal systems of most European languages and even to Belgian French, which uses transparent forms like \emph{nonante-sept} (ninety-seven).
How do large language models (LLMs) handle this linguistic irregularity?
Do they simply memorize the mapping, or does the vigesimal structure leave traces in their internal representations?

Understanding how LLMs bridge linguistic form and mathematical meaning is a central question in the study of numerical reasoning.
Recent work has shown that models encode numbers using base-10 digit representations \citep{levy2024language}, that sinusoidal probes can recover numeric values from hidden states \citep{kadlcik2025pretrained, stefanik2025unravelling}, and that implicit compositional operators embedded in numeral systems pose a significant challenge \citep{bhattacharya2025investigating}.
However, no prior study has examined how the French vigesimal system---one of the most widely encountered irregular number systems---is represented internally by LLMs.

We address this gap with the first systematic study of French vigesimal number processing in LLMs.
Our approach combines behavioral evaluation using \gptmodel with representational analysis of \mistral's hidden states, leveraging Belgian French as a natural control condition that isolates the effect of vigesimal structure from other linguistic factors.

Our behavioral results show that \gptmodel achieves perfect accuracy (100\%) on all French number tasks, including vigesimal conversion, arithmetic, and counting across vigesimal boundaries.
Yet beneath this behavioral ceiling, \mistral's internal representations tell a different story: vigesimal numbers cluster by linguistic base rather than numeric proximity, with 70\% of nearest neighbors sharing a prefix like \quatrevingt.
French number words yield higher linear probe accuracy ($R^2 = 0.979$) than digit strings ($R^2 = 0.943$), and Belgian French produces the most numerically ordered representations ($\rho = 0.591$).

We make the following contributions:
\begin{itemize}[leftmargin=*,itemsep=0pt,topsep=0pt]
    \item We conduct the first behavioral and representational analysis of French vigesimal number processing in LLMs, demonstrating perfect behavioral accuracy alongside structured internal representations.
    \item We show that vigesimal numbers cluster by linguistic base in representation space, revealing that LLMs organize these numbers by their surface form rather than their numeric value.
    \item We establish Belgian French as an effective control condition, finding that its regular decimal system produces representations with superior ordinal structure compared to standard French.
    \item We find that French number words produce more linearly decodable numeric representations than raw digit strings, suggesting that linguistic context enriches rather than obscures numerical information.
\end{itemize}
